\clearpage
\fancyhead[L]{\small GAUSS AN BODE.}
\hypertarget{native-444-part-lower}{}
\begin{center}
{\large GAUSS AN BODE.}

\noindent\rule{\linewidth}{0.4pt}

Astronomisches Jahrbuch für 1825. Seite 103..105. Berlin 1821.

\noindent\rule{\linewidth}{0.4pt}
\end{center}

\begin{flushright}
Göttingen, den 26. Dec. 1821.
\end{flushright}

Sie werden mich wegen meines so langen Stillschweigens entschuldigt halten, wenn ich Ihnen
sage, dass ich den grössten Theil des Jahres von hier abwesend gewesen bin und selbst noch im Spätherbst eine Reise nach Altona zur Empfangnahme des \textsc{Ramsden}schen Zenithsectors gemacht habe, von
wo ich erst vor kurzem hieher zurückgekommen bin.

Bei meiner Triangulation, wo ich bisher an fünf Dreieckspunkten die Winkel gemessen habe,
habe ich die Dreiecke so gross wie möglich zu machen gesucht. Ueber das neue, von mir zu diesem
Behuf angewandte Hülfsmittel, den \emph{Heliotrop}, und die ersten damit gemachten, ins Grosse gehenden
Versuche, werden Sie die Nachricht in Nr. 126. der hiesigen gelehrten Anzeigen [B.\,IV.\,d.\,W.] gelesen
haben. Seit der Zeit habe ich davon beständig Gebrauch gemacht, nicht allein als Zielpunkt beim Winkelmessen, sondern auch mit nicht weniger glücklichem Erfolg zu telegraphischen Signalisirungen. Die
gewaltige Wirkung des reflectirten Sonnenlichts von einem Spiegel von 2 Zoll Breite und $1\tfrac{1}{4}$ Zoll Höhe,
welches in Entfernungen von 5, 6, $7\tfrac{1}{2}$, ja einmal von $9\tfrac{1}{3}$ geographischen Meilen mit \emph{blossen Augen} gesehen wurde, pflegt diejenigen, die sie zum ersten male erfahren, und nicht durch theoretische Berechnung darauf vorbereitet sind, gewöhnlich in Erstaunen zu setzen. Bei einem nur einigermassen günstigen Zustande der Luft giebt es jetzt für die Grösse der Dreiecksseiten keine Grenzen mehr, als die
die Krümmung der Erde setzt, zumal wenn man, wie ich es bei zwei neu angefertigten Heliotropen von
ganz verschiedener Construction gethan habe, den Spiegeln noch etwas grössere Dimensionen giebt.

Die erwähnten Beschäftigungen haben mich bis jetzt gehindert, die Berechnung einer Ephemeride
für die nächste Erscheinung der Pallas zu besorgen und die früheren von mir hier gemachten Beobachtungen für Ihr Jahrbuch in Ordnung zu bringen. Nur meine Beobachtungen des Cometen von
diesem Jahre kann ich Ihnen jetzt hier beifügen.

\Needspace{15\baselineskip}
\begin{center}
\small
\setlength{\tabcolsep}{4pt}
\begin{tabular}{l|rrr|rrr|rrr}
1821 & \multicolumn{3}{c|}{\begin{tabular}{c}Mittlere Zeit\\ in Göttingen\end{tabular}} & \multicolumn{3}{c|}{Gerade Aufsteigung} & \multicolumn{3}{c}{Nordliche Abweichung} \\
\hline
Jan. 30 & 7\textsuperscript{u} & 34\textsuperscript{m} & 42\textsuperscript{s} & $359^\circ$ & $27'$ & $7''$ & $16^\circ$ & $4'$ & $36''$ \\
Febr. 3 & 7 & 3 & 56 & 359 & 3 & 54 & 15 & 46 & 3 \\
 7 & 6 & 42 & 38 & 358 & 45 & 5 & 15 & 29 & 49 \\
 9 & 6 & 42 & 30 & 358 & 36 & 24 & 15 & 21 & 22 \\
 10 & 6 & 52 & 27 & 358 & 32 & 19 & 16 & 17 & 50 \\
 11 & 7 & 12 & 0 & 358 & 28 & 27 & 15 & 14 & 26 \\
März 1 & 7 & 18 & 7 & 357 & 13 & 34 & 14 & 8 & 37 \\
 5 & 7 & 5 & 28 & 356 & 54 & 11 & 13 & 43 & 5 \\
\end{tabular}
\end{center}

\Needspace{18\baselineskip}
\SourceLeaf{445}
Die parabolischen Elemente, velche Hr. von \textsc{Staudt} hiernach berechnet hat, sind folgende:

\begin{center}
\small
\begin{tabular}{p{.43\linewidth}p{.49\linewidth}}
Sonnennähe Zeit \mbox{.\,.\,.\,.\,.\,.\,.} & 1821 März $21.61890$ in Göttingen \\
—\quad Argum. der Breite \mbox{.\,.\,.\,.\,.\,.} & $169^\circ\;10'\;9''3$ \\
—\quad Logarithm des Abstandes \mbox{.\,.\,.} & $8.9645990$ \\
{\rightskip=0pt plus 3em Länge des aufsteigenden Knoten \mbox{.\,.\,.\,.\,.}\par} & $48^\circ\;44'\;14''7$ vom m. Aequin. d. 1. Febr. \\
Neigung der Bahn \mbox{.\,.\,.\,.\,.\,.\,.\,.\,.\,.} & $106\quad40\quad16.4$ \\
\end{tabular}
\end{center}

Herr von \textsc{Staudt} hat diese Elemente nicht bloss mit meinen Beobachtungen, sondern auch mit allen andern damals bekannt gewordenen verglichen; die Resultate dieser Vergleichung lege ich hier für Ihr Jahrbuch von seiner eigenen Hand bei, da die Zeit zu kurz ist, sie abzuschreiben.

\Needspace{34\baselineskip}
\begin{center}
\begin{tabular}{l|r|r}
& \multicolumn{2}{c}{Differenzen in} \\
& AR. & Decl. \\
\multicolumn{3}{c}{Paris} \\
Jan. 21 & $+0''2$ & $+0''1$ \\
\multicolumn{3}{c}{Göttingen} \\
Jan. 30 & $-18''8$ & $+34''9$ \\
Febr. 3 & $+3.9$ & $-6.0$ \\
7 & $-13.1$ & $-36.0\!:$ \\
9 & $-0.1$ & $+14.7$ \\
10 & $+2.6$ & $+7.3$ \\
11 & $-1.2$ & $-3.5$ \\
März 1 & $-0.3$ & $-0.3$ \\
5 & $-5.3$ & $-20.6$ \\
\multicolumn{3}{c}{Bremen} \\
Jan. 30 & $-12.6$ & $+12.4$ \\
30 & $+9.1$ & $+34.2$ \\
Febr. 2 & $+24.7$ & $+9.7$ \\
5 & $-7.3$ & $-42.3\!::$ \\
7 & $+8.8$ & $+16.0\!:$ \\
8 & $+4.9$ & $+22.4$ \\
9 & $+5.0$ & $+14.1$ \\
10 & $-5.8$ & $+20.2$ \\
11 & $+4.4$ & $+3.2$ \\
12 & $-9.5$ & $+0.4$ \\
13 & $+0.6$ & $-24.6$ \\
14 & $-2.8$ & $-19.6$ \\
19 & $+10.2$ & $-3.4$ \\
März 1 & $+8.5$ & $-9.4$ \\
5 & $-0.6$ & $-7.6$ \\
6 & $+8.6$ & $-0.2$ \\
\end{tabular}
\end{center}

\Needspace{38\baselineskip}
\begin{center}
\begin{tabular}{l|r|r}
& \multicolumn{2}{c}{Differenzen in} \\
& AR. & Decl. \\
\multicolumn{3}{c}{Mannheim} \\
Febr. 6 & $-0''2$ & $+28''9\!:$ \\
7 & $+2.2$ & $+1.3$ \\
8 & $+9.6$ & $-1.1$ \\
9 & $-6.2$ & $+5.7$ \\
10 & $+5.9$ & $+7.3$ \\
11 & $+12.0$ & $+17.4$ \\
12 & $+8.6$ & $+6.3$ \\
13 & $+8.6$ & $+4.2$ \\
14 & $+9.2$ & $+12.7$ \\
15 & $+9.4$ & $-3.7$ \\
27 & $+3.7$ & $-18.2$ \\
\multicolumn{3}{c}{Königsberg} \\
Febr. 9 & $-1.2$ & $+0.5$ \\
10 & $-5.4$ & $+2.9$ \\
11 & $+1.2$ & $-7.6$ \\
12 & $+1.9$ & $-1.1$ \\
14 & $+14.3$ & $-6.3$ \\
15 & $+11.7$ & $-20.8$ \\
19 & $-9.5$ & $+20.7$ \\
25 & $+35.3$ & $+35.9$ \\
27 & $+5.1$ & $-10.6$ \\
März 4 & $+15.6$ & $-12.4$ \\
5 & $-11.4$ & $+4.3$ \\
6 & $+7.8$ & $-17.2$ \\
\multicolumn{3}{c}{Hamburg} \\
Febr. 7 & $-3.2$ & $-29.2$ \\
8 & $+4.5$ & $+28.9$ \\
9 & $-6.5$ & $+9.9$ \\
10 & $+4.5$ & \mbox{.\,.\,.\,.} \\
12 & $-7.0$ & \mbox{.\,.\,.\,.} \\
\multicolumn{3}{c}{Dorpat} \\
Febr. 14 & $-3.6$ & $+83.0\!::$ \\
\end{tabular}
\end{center}

\Needspace{39\baselineskip}
\begin{center}
\begin{tabular}{l|r|r}
& \multicolumn{2}{c}{Differenzen in} \\
& AR. & Decl. \\
\multicolumn{3}{c}{Seeberg} \\
Febr. 3 & $+30''2$ & $+103''4\!::$ \\
5 & $-8.6$ & $-41.5$ \\
7 & $+14.3$ & $-95.1\!::$ \\
8 & $+25.1$ & $+40.3\!::$ \\
9 & $+9.8$ & $-20.6$ \\
10 & $-4.9$ & $-19.0$ \\
11 & $+0.1$ & $-15.1$ \\
12 & $+10.2$ & $-20.0$ \\
14 & $+5.3$ & $-39.7$ \\
\multicolumn{3}{c}{Wien} \\
Febr. 20 & $-78.8$ & $-28.1$ \\
\multicolumn{3}{c}{Mailand} \\
Jan. 31 & $-37.7$ & $+4.5$ \\
Febr. 1 & $-22.9$ & $-4.4$ \\
2 & $-0.2$ & $-3.1$ \\
3 & $-29.4$ & $-0.2$ \\
7 & $-31.5$ & $-10.5$ \\
8 & $-19.8$ & $-2.3$ \\
\multicolumn{3}{c}{Padua} \\
Febr. 16 & $-50.5$ & $-56.1$ \\
16 & $-30.8$ & $-60.2$ \\
17 & $-69.4$ & $-15.5$ \\
17 & $-44.0$ & $-45.9$ \\
18 & $-61.5$ & $-1.3$ \\
18 & $-34.6$ & $-30.1$ \\
\multicolumn{3}{c}{Florenz} \\
Febr. 16 & $-14.8$ & $-24.9$ \\
17 & $-7.0$ & $-12.9$ \\
21 & $+5.9$ & $-47.6$ \\
22 & $-19.1$ & $-34.8$ \\
23 & $-8.6$ & $-47.5$ \\
25 & $-26.3$ & $-265.0$ \\
27 & $+16.0$ & $-603.2$ \\
\end{tabular}
\end{center}
